\section{Introduction}\label{sec:intro}
Many engineering decisions rest on a simulation that is too slow to run as often as one would like. Choosing the shape of a part, finding how uncertain an input leaves the answer, or fitting a model to measurements can each need thousands of runs. A common remedy is a \emph{surrogate}, or stand-in model: the simulation is run a limited number of times, say a few hundred, a cheap model is trained on the inputs and outputs, and the cheap model then replaces the simulation in all further work~\cite{quarteroni2016}. For example, if one run takes an hour, five hundred runs take about three weeks on one machine, whereas a trained stand-in answers in milliseconds.

The practical question is which kind of stand-in to build, and how many runs are enough. One long-standing route compresses the output fields with proper orthogonal decomposition (POD) and regresses the weights of the leading modes on the inputs, with a neural network~\cite{hesthaven2018,bhattacharya2021} or a Gaussian process~\cite{guo2018}. The other route trains a neural network to produce the whole field~\cite{fresca2022,franco2023}. Published comparisons are mixed. A survey of machine-learning solvers for fluid problems found that most papers claiming to beat a standard method used a weak baseline~\cite{mcgreivy2024}; kernel and Gaussian-process methods have been reported to be competitive with DeepONet~\cite{lu2021} and the Fourier neural operator~\cite{li2021} on benchmark problems~\cite{batlle2024}; and no universal winner was found between neural and polynomial surrogates~\cite{westermann2026}. A practitioner holding a few hundred runs of their own simulation still has no procedure that tells them which stand-in to use on \emph{their} data, how far it can be trusted, or how many more runs it would take to reach a target accuracy.

\paragraph{Gap} Existing comparisons rank methods on fixed problem sets and do not give the practitioner a validated procedure to select a stand-in, to see whether the POD basis limits accuracy, and to forecast the number of runs needed.

\subsection*{Motivation and contributions}
\begin{enumerate}
\item \textbf{A selection tool, SurroSelect} (\cref{sec:method}). It fits a pool of stand-in models, chooses by cross-validation while reporting which models are tied with the best, checks whether the POD basis is the bottleneck, and forecasts the error for more runs from a learning curve. It is free software with a command-line interface (Sec.~\ref{sec:S-tool}).
\item \textbf{A pre-specified evaluation on problems solved by third-party code} (\cref{sec:results-tool}). Seven criteria were committed to version control before the evaluation runs on eight problems solved by the third-party finite-element library \texttt{scikit-fem} (smoke tests of the tool on three of them preceded the commit; see \cref{sec:discussion}, limitation 2). We report which criteria were met and which were not, against a neural network tuned with a larger budget than any other method, and on two real measured datasets.
\item \textbf{A \NAllProblems-problem study of where each family wins} (\cref{sec:results-study,sec:results-mech}). Controls separate the effects of the reduced basis, of the regressor, and of how widely its hyper-parameters are searched, and the outcome of a comparison is shown to depend on the training budget and the tuning effort.
\item \textbf{A record of what did not work} (\cref{sec:discussion}), including a pilot-data diagnostic that was not confirmed, a control that had to be widened after review, and every deviation from the plan (Supplement, Sec.~\ref{sec:S-deviations}).
\end{enumerate}
The paper proposes no new regression model; the pool consists of classical methods, and the contribution is the procedure, its validation, and the evidence about when each member wins.

\subsection*{Organisation}
\Cref{sec:related} reviews related work; \cref{sec:prelim} fixes notation and assumptions; \cref{sec:method} describes the method; \cref{sec:setup} the experimental framework; \cref{sec:results} the results; \cref{sec:discussion} the discussion and limitations; \cref{sec:conclusion} concludes.
